\documentclass[11pt,a4paper]{article}
\usepackage[margin=2.5cm]{geometry}
\usepackage{amsmath,amssymb,bm}
\usepackage{microtype}
\usepackage[hidelinks]{hyperref}
\newcommand{\vect}[1]{\bm{#1}}
\newcommand{\dd}{\mathop{}\!\mathrm{d}}
\newcommand{\erf}{\operatorname{erf}}
\newcommand{\ee}{\mathrm{e}}

\title{Theory Notes: Gaussian-Surface COSMO/C-PCM Coupled to GFN2-xTB}
\author{Working notes for the tblite implementation}
\date{\today}

\begin{document}
\maketitle

\section{Purpose and scope}

The implementation couples GFN2-xTB in \texttt{tblite} to a conductor-like
continuum model using the iSwiG cavity from \texttt{moist}.  The solute can be
represented either by the complete GFN2 AO density matrix or by atom-centred
GFN2 monopoles and optional dipoles and quadrupoles.  Both descriptions interact
with exactly the same Gaussian functions on the cavity.  Thus, the intended
variable in a comparison is the solute-density representation, not the cavity
or surface-charge representation.

\section{iSwiG cavity}

The cavity is a union of atom-centred spheres.  The default radii are obtained
from \texttt{moist}'s \texttt{new\_gauss\_radii} model.  This table contains
Bondi radii where available, uses $R_{\mathrm H}=1.10$~\AA, uses Mantina
\emph{et al.} radii for missing main-group elements, and uses $2.00$~\AA{} as
a fallback.  All base radii are multiplied by 1.2.  For example,
\begin{equation}
 R_{\mathrm H}^{\mathrm{cav}}=1.2(1.10~\text{\AA})=1.32~\text{\AA}.
\end{equation}
Explicit custom radii remain possible.

The current default is a fixed 110-point Lebedev grid on every sphere.  This
value is passed explicitly to \texttt{moist}; its internal grid default is not
used.  The iSwiG switching function smoothly suppresses buried points.  For
each retained point $i$, \texttt{moist} provides its position $\vect{s}_i$,
area $a_i$, Lebedev weight $w_i$, switching value $F_i$, Gaussian width
$\xi_i$, and owner atom.

The normalized surface Gaussian is
\begin{equation}
 g_i(\vect{r})=
 \left(\frac{\xi_i^2}{\pi}\right)^{3/2}
 \exp[-\xi_i^2|\vect{r}-\vect{s}_i|^2],
 \qquad \int g_i(\vect{r})\dd\vect{r}=1,
 \label{eq:gi}
\end{equation}
and the apparent surface density is
\begin{equation}
 \rho_{\mathrm{surf}}(\vect{r})=\sum_i q_i g_i(\vect{r}).
\end{equation}

\section{Gaussian surface interaction and dielectric equation}

The surface matrix is the Coulomb metric
\begin{equation}
 A_{ij}=(g_i|g_j).
\end{equation}
The entries constructed by \texttt{moist} are
\begin{align}
 A_{ij}&=\frac{\erf(\xi_{ij}r_{ij})}{r_{ij}},\quad i\ne j,\\
 \xi_{ij}&=\frac{\xi_i\xi_j}{\sqrt{\xi_i^2+\xi_j^2}},\\
 A_{ii}&=\sqrt{\frac{2}{\pi}}\frac{\xi_i}{F_i}.
\end{align}
This matrix definition has been tested in \texttt{moist} against ORCA 6.1.1
surface charges and dielectric energy for a supplied ORCA surface.

For the solute potential vector $\vect{\phi}$, the stationary functional is
\begin{equation}
 \mathcal G_{\mathrm{pol}}(\vect q;\vect\phi)
 =\frac{1}{2f_\epsilon}\vect q^{\mathrm T}A\vect q
  +\vect q^{\mathrm T}\vect\phi.
\end{equation}
Stationarity gives
\begin{equation}
 A\vect q=-f_\epsilon\vect\phi,
\end{equation}
where
\begin{align}
 f_\epsilon^{\mathrm{COSMO}}&=\frac{\epsilon-1}{\epsilon+1/2},\\
 f_\epsilon^{\mathrm{CPCM}}&=\frac{\epsilon-1}{\epsilon}.
\end{align}
The polarization energy is
\begin{equation}
 E_{\mathrm{pol}}=\frac12\vect q^{\mathrm T}\vect\phi
 =-\frac{f_\epsilon}{2}\vect\phi^{\mathrm T}A^{-1}\vect\phi.
 \label{eq:epol}
\end{equation}

\section{Full GFN2-xTB density}

GFN2-xTB uses a minimal valence AO basis $\{\chi_\mu\}$.  The charge-density
matrix is the sum of the spin blocks,
\begin{equation}
 P_{\mu\nu}=\sum_\sigma P_{\mu\nu}^{\sigma},
\end{equation}
and the electronic number density is
\begin{equation}
 \rho_{\mathrm e}(\vect r)=
 \sum_{\mu\nu}P_{\mu\nu}\chi_\mu(\vect r)\chi_\nu(\vect r).
\end{equation}
``Full density'' therefore means the complete density within the GFN2 valence
AO model; it does not mean a conventional all-electron density in a large basis.
The nuclei and excluded core electrons are represented by xTB effective core
charges $Z_A^{\mathrm{eff}}$, stored as \texttt{wfn\%n0at}.  The total solute
charge density is
\begin{equation}
 \rho_{\mathrm{solute}}(\vect r)=
 \sum_A Z_A^{\mathrm{eff}}\delta(\vect r-\vect R_A)
 -\rho_{\mathrm e}(\vect r).
\end{equation}

The effective-core interaction with surface Gaussian $i$ is
\begin{equation}
 U_{iA}=\frac{\erf(\xi_iR_{iA})}{R_{iA}},
 \qquad R_{iA}=|\vect s_i-\vect R_A|.
\end{equation}
The AO-pair interaction is the three-centre two-electron integral
\begin{equation}
 B_{i,\mu\nu}=(g_i|\chi_\mu\chi_\nu)
 =\iint\frac{g_i(\vect s)\chi_\mu(\vect r)\chi_\nu(\vect r)}
 {|\vect s-\vect r|}\dd\vect s\dd\vect r,
\end{equation}
which is evaluated with \texttt{libcint}.  Hence
\begin{equation}
 \phi_i^{\mathrm{full}}
 =\sum_A Z_A^{\mathrm{eff}}U_{iA}
 -\sum_{\mu\nu}P_{\mu\nu}B_{i,\mu\nu}.
 \label{eq:fullphi}
\end{equation}
The stationary reaction-field AO potential returned to the SCF is
\begin{equation}
 V_{\mu\nu}^{\mathrm{reac}}=-\sum_i q_iB_{i,\mu\nu}.
\end{equation}
There is no factor $1/2$ in this potential: differentiating the stationary
energy in Eq.~\eqref{eq:epol} supplies the complete response.

\section{Multipole representation}

The alternative representation uses the converged atom-resolved GFN2 variables
$q_A$, $\vect\mu_A$, and $Q_A$.  The monopoles are always used.  Dipoles and
quadrupoles are independently enabled.  The Cartesian quadrupole ordering is
\begin{equation}
 (Q_{xx},Q_{xy},Q_{yy},Q_{xz},Q_{yz},Q_{zz}),
 \qquad \operatorname{Tr}Q=0.
\end{equation}

\subsection{Explicit derivation of the Gaussian--multipole coupling}

This subsection records how the implemented coupling was derived.  It was not
copied from an ORCA or \texttt{moist} routine.  It follows by treating point
multipoles as derivatives of a delta distribution and differentiating the
exact Gaussian--point-charge interaction.

Let
\begin{equation}
 \vect R=\vect s_i-\vect R_A,
 \qquad R=|\vect R|,
 \qquad x=\xi_iR,
\end{equation}
and define
\begin{equation}
 U(R)=(g_i|\delta_A)=\frac{\erf(\xi_iR)}{R}.
 \label{eq:screened-U}
\end{equation}
For reference, the point-charge limit is recovered when $x\to\infty$:
\begin{equation}
 U(R)\longrightarrow\frac1R.
\end{equation}

\subsubsection*{Monopole}

A monopole charge distribution is
\begin{equation}
 \rho_A^{(0)}(\vect r)=q_A\delta(\vect r-\vect R_A),
\end{equation}
so its contribution is simply
\begin{equation}
 \phi_{iA}^{(0)}=q_AU(R).
\end{equation}

\subsubsection*{Dipole}

With the sign convention used here, a point dipole is generated by a derivative
with respect to its source centre,
\begin{equation}
 \rho_A^{(1)}(\vect r)
 =\vect\mu_A\mathbin{\cdot}\nabla_{R_A}
   \delta(\vect r-\vect R_A).
\end{equation}
Because $\vect R=\vect s_i-\vect R_A$,
\begin{equation}
 \nabla_{R_A}U(R)=-\nabla_RU(R)=-\frac{U'(R)}{R}\vect R.
 \label{eq:source-gradient}
\end{equation}
Differentiating Eq.~\eqref{eq:screened-U} gives
\begin{equation}
 U'(R)=\frac{2\xi_i}{\sqrt\pi}\frac{\ee^{-x^2}}{R}
       -\frac{\erf(x)}{R^2}.
\end{equation}
Inserting this into Eq.~\eqref{eq:source-gradient} yields
\begin{align}
 \vect C_{iA}^{(1)}&=K_1(R)\vect R,\\
 K_1(R)&=
 \frac{\erf(x)}{R^3}
 -\frac{2\xi_i\ee^{-x^2}}{\sqrt\pi R^2}.
 \label{eq:k1}
\end{align}
Thus
\begin{equation}
 \phi_{iA}^{(1)}=\vect\mu_A\mathbin{\cdot}\vect C_{iA}^{(1)}.
\end{equation}
As $x\to\infty$, the exponential vanishes and
\begin{equation}
 \vect C_{iA}^{(1)}\longrightarrow\frac{\vect R}{R^3},
\end{equation}
which is exactly the former \texttt{tblite} point-dipole kernel.  This
asymptotic check fixes the sign convention.

\subsubsection*{Quadrupole}

For any radial function $U(R)$, its Cartesian Hessian is
\begin{equation}
 \frac{\partial^2U}{\partial R_\alpha\partial R_\beta}
 =\frac{U'(R)}{R}\delta_{\alpha\beta}
 +\left(\frac{U''(R)}{R^2}-\frac{U'(R)}{R^3}\right)
 R_\alpha R_\beta.
 \label{eq:radial-hessian}
\end{equation}
The existing \texttt{tblite} point-quadrupole potential uses
\begin{equation}
 Q_{\alpha\beta}\frac{R_\alpha R_\beta}{R^5}.
 \label{eq:tblite-point-q}
\end{equation}
This is related to the Hessian of $1/R$ by
\begin{equation}
 \frac13 Q_{\alpha\beta}
 \frac{\partial^2}{\partial R_\alpha\partial R_\beta}\frac1R
 =Q_{\alpha\beta}\frac{R_\alpha R_\beta}{R^5},
 \label{eq:q-factor-third}
\end{equation}
because
\begin{equation}
 \frac{\partial^2}{\partial R_\alpha\partial R_\beta}\frac1R
 =\frac{3R_\alpha R_\beta-R^2\delta_{\alpha\beta}}{R^5}
\end{equation}
and $Q_{\alpha\alpha}=0$.  Equation~\eqref{eq:q-factor-third} explains both
the factor $1/3$ and why the isotropic part of Eq.~\eqref{eq:radial-hessian}
can be omitted for the traceless GFN2 quadrupoles.

The Gaussian quadrupole interaction is therefore defined by applying the same
operator to $U(R)$:
\begin{equation}
 \phi_{iA}^{(2)}=
 \frac13Q_{A,\alpha\beta}
 \frac{\partial^2U(R)}{\partial R_\alpha\partial R_\beta}.
 \label{eq:q-gaussian-definition}
\end{equation}
The required second derivative is
\begin{equation}
 U''(R)=
 -\frac{4\xi_i^3}{\sqrt\pi}\ee^{-x^2}
 -\frac{4\xi_i}{\sqrt\pi R^2}\ee^{-x^2}
 +\frac{2\erf(x)}{R^3}.
\end{equation}
Using Eqs.~\eqref{eq:radial-hessian} and
\eqref{eq:q-gaussian-definition}, the radial factor multiplying
$Q_{\alpha\beta}R_\alpha R_\beta$ is
\begin{align}
 K_2(R)
 &=\frac13\left(\frac{U''(R)}{R^2}
                 -\frac{U'(R)}{R^3}\right)\\
 &=\frac{\erf(x)}{R^5}
 -\frac{2\xi_i\ee^{-x^2}}{\sqrt\pi R^4}
 -\frac{4\xi_i^3\ee^{-x^2}}{3\sqrt\pi R^2}.
 \label{eq:k2}
\end{align}
In the six-component storage convention, the implemented kernel is
\begin{equation}
 \vect C_{iA}^{(2)}=K_2(R)
 \begin{pmatrix}
 X^2 & 2XY & Y^2 & 2XZ & 2YZ & Z^2
 \end{pmatrix}^{\mathrm T}.
\end{equation}
The factors of two are the two equal off-diagonal terms in the full Cartesian
contraction.  In the point-charge limit,
\begin{equation}
 K_2(R)\longrightarrow\frac1{R^5},
\end{equation}
so the previous point-quadrupole kernel is recovered exactly.

\subsubsection*{Combined multipole potential}

The surface potential from the atom-centred representation is
\begin{equation}
 \phi_i^{\mathrm{mult}}=\sum_A\left[
 q_AC_{iA}^{(0)}
 +\vect\mu_A\mathbin{\cdot}\vect C_{iA}^{(1)}
 +\vect Q_A\mathbin{\cdot}\vect C_{iA}^{(2)}
 \right],
 \label{eq:multipole-potential}
\end{equation}
with optional omission of the dipole or quadrupole term.  The same coupling
matrices are used to form $\vect\phi$ and to return the adjoint reaction-field
potentials to the GFN2 SCF variables.

\paragraph{Validation status.}
The monopole result is the standard Gaussian--point-charge integral, and the
dipole and quadrupole expressions follow analytically by differentiation.  The
signs and normalization have also been checked against the original point
kernels in the $x\to\infty$ limit.  Nevertheless, before publication the
quadrupole convention should additionally be tested against direct numerical
integration and finite differences for several Gaussian widths and geometries.

\section{Controlled comparison of density representations}

The full-density and multipole calculations use the same geometry, radii,
Lebedev grid, point filtering, surface Gaussians, switching values, matrix $A$,
dielectric factor, linear solver, and energy expression.  Only
Eq.~\eqref{eq:fullphi} versus Eq.~\eqref{eq:multipole-potential} changes.
Useful calculations are monopoles only; monopoles plus dipoles; monopoles plus
quadrupoles; all multipoles through quadrupole order; and the complete AO
density.  Potential errors can be quantified through
\begin{align}
 \Delta_{\mathrm{RMS}}(\phi)&=
 \sqrt{\frac1N\sum_i(\phi_i^{\mathrm{mult}}-
 \phi_i^{\mathrm{full}})^2},\\
 \Delta_{\max}(\phi)&=\max_i|\phi_i^{\mathrm{mult}}-
 \phi_i^{\mathrm{full}}|.
\end{align}

\section{ORCA-style output and present limitations}

The file \texttt{tblite.cpcm} is overwritten during every energy evaluation,
leaving the final SCF surface after convergence.  It contains coordinates,
area, potential, conductor charge, Lebedev weight, switching value, Gaussian
width, and zero-based owner atom.  The printed conductor charge is
$q_i/f_\epsilon$; internally, $q_i$ includes dielectric scaling.

Important limitations and checks are:
\begin{enumerate}
 \item the fixed 110-point grid differs from ORCA's recent adaptive default;
 \item the AO density is complete only within the GFN2 valence basis;
 \item the three-centre surface integrals require continuing independent tests;
 \item the Gaussian quadrupole convention requires direct numerical tests;
 \item a complete analytic nuclear gradient is not yet implemented;
 \item the model described here contains only the electrostatic continuum term.
\end{enumerate}

\begin{thebibliography}{9}
\bibitem{bondi} A. Bondi, \emph{J. Phys. Chem.} \textbf{68}, 441--451 (1964).
\bibitem{mantina} M. Mantina \emph{et al.}, \emph{J. Phys. Chem. A}
\textbf{113}, 5806--5812 (2009).
\bibitem{iswig} A. W. Lange and J. M. Herbert, \emph{J. Chem. Phys.}
\textbf{133}, 244111 (2010).
\bibitem{gfn2} C. Bannwarth, S. Ehlert, and S. Grimme,
\emph{J. Chem. Theory Comput.} \textbf{15}, 1652--1671 (2019).
\end{thebibliography}

\end{document}
